\documentclass[../../main]{subfiles}
\begin{document}
\part{Category theory}

\ifSubfilesClassLoaded{%
    \tableofcontents
}{}

\chapter*{Introduction}

The development of sheaf cohomology in the mid-20th century marked a fundamental shift in how we understand the relationship between local and global data in mathematics.
At the center of this revolution was the language of Category Theory.
Rather than focusing solely on the internal elements of a set or the points of a space, category theory prioritizes the relationships—the morphisms—between objects.

Historically, the seeds of this language were sown in the 1940s. Samuel Eilenberg and Saunders Mac Lane initially introduced the concepts of ``categories'', ``functors'', and "natural transformations" in their 1945 paper, General Theory of Natural Equivalences, which was first rejected on the ground that it had no content. Interestingly, their primary motivation was not to create a new branch of mathematics, but to provide a formal definition of a "natural transformation" to compare different ways of calculating homology groups in algebraic topology.

As the 1950s progressed, the utility of this language exploded. Alexander Grothendieck utilized categorical foundations to revolutionize algebraic geometry, moving the focus from specific varieties to the "sheaves" of functions living on them. By treating sheaves as objects within a category, Grothendieck and his school (notably in \cite{SGA4}) were able to define cohomology not just as a topological invariant, but as a "derived functor"—a universal construction that measures the failure of an operation to preserve exactness.  

The objective of this first part is to establish the rigorous categorical vocabulary that emerged from this period, which now serves as the "operating system" for modern homological algebra.

\begin{enumerate}
  \item \hyperref[chap:categories-morphisms-duality]{Categories, Morphisms, and Duality} introduces the basic arena of our study. We define the category of open sets \(\Open(X)\), which acts as the foundational "site" for any sheaf construction.  

  \item \hyperref[chap:functors-natural-transformations]{Functors and Natural Transformations} defines the "mappings between categories". This allows us to treat a sheaf as a contravariant functor and provides the mechanism for "sheafification" via natural transformations.  

  \item \hyperref[chap:limits-colimits]{Limits and Colimits} provides the tools for "categorical calculus". These universal constructions are vital for defining stalks as direct limits and the gluing property of sheaves as equalizers.  

  \item \hyperref[chap:more-functors]{More on Functors} explores the deeper symmetries of the theory through Adjunctions and the Yoneda Lemma. These results provide the justification for why the global sections functor \(\Gamma(X, -)\) is left-exact, directly motivating the eventual construction of right derived functors.
        
\end{enumerate}

By the end of this part, we will have constructed the scaffolding necessary to define Abelian Categories, the specific environment where homological algebra and sheaf cohomology come to life.

The main references used for this part are \cite{riehlCategoryTheoryContext}, \cite{aluffiAlgebraChapter02009} and \cite{maclaneCategoriesWorkingMathematician1978}.

\subfile{1.Categories-morphisms-and-duality}

\subfile{2.Functors-and-natural-transformations}

\subfile{3.Limits-and-colimits}

\subfile{4.More-on-functors}

\end{document}
